%
% architecture.tex
%
% Dungeon: Technical Architecture Specification
% Written for engineers implementing the system.
%
% Compile: pdflatex architecture.tex
%

\documentclass[a4paper,10pt]{article}
\usepackage[margin=1in]{geometry}
\usepackage{booktabs}
\usepackage{listings}
\usepackage{xcolor}
\usepackage{enumitem}
\usepackage{longtable}
\usepackage{textcomp}
\usepackage[colorlinks=true,linkcolor=blue,citecolor=blue,urlcolor=blue]{hyperref}

\lstset{
  basicstyle=\ttfamily\small,
  keywordstyle=\color{blue!80!black},
  commentstyle=\color{green!50!black},
  stringstyle=\color{red!60!black},
  breaklines=true,
  frame=single,
  framerule=0.3pt,
  rulecolor=\color{gray!50},
  backgroundcolor=\color{gray!5},
  columns=fullflexible,
}

\newcommand{\iface}[1]{\texttt{#1}}
\newcommand{\tool}[1]{\texttt{#1}}
\newcommand{\cmd}[1]{\texttt{#1}}
\newcommand{\field}[1]{\texttt{#1}}
\newcommand{\mode}[1]{\textbf{#1}}
\newcommand{\pkg}[1]{\texttt{#1}}

\begin{document}

\title{Crypt: Technical Architecture Specification}
\author{LLM-Directed Text Adventure with Go Game Engine}
\date{March 2026 --- v0.4}
\hypersetup{
  pdftitle={Crypt: Technical Architecture Specification},
  pdfauthor={Punt Labs},
  pdfsubject={Architecture},
  pdfcreator={pdflatex}
}
\maketitle

\tableofcontents
\newpage

% ===================================================================
\section{Overview}
% ===================================================================

Crypt is a text adventure game for Claude Code.  In its current
form the entire engine --- parser, rules, narrator, state machine ---
is a single LLM skill prompt.  This specification describes the
evolved architecture: a deterministic Go game engine that handles
all state transitions, with the LLM acting as Dungeon Master
(narrator, cartographer, encounter designer) rather than parser.

The central insight is \textbf{play mode as a first-class concept}.
A play mode is a named combination of three swappable Go interfaces:
\iface{CommandInterpreter}, \iface{Narrator}, and \iface{Renderer}.
All modes share the same engine, the same save files, and the same
scenario format.  Switching modes mid-adventure --- save in
\mode{dm}, resume in \mode{solo} --- is explicitly supported.

\paragraph{Note on companion documents.}
Three binaries now exist: \cmd{cryptd} (server), \cmd{crypt} (client),
and \cmd{crypt-admin} (scenario authoring).  The authoritative reference
for implemented game mechanics (combat formulas, damage reduction, leveling
tables, point-buy stats, consumables) is \texttt{docs/gameplay.tex}.
This architecture document describes the system-level design; some
mechanics sections below reflect design-stage aspirations that differ
from the current implementation.

\subsection{Mechanics Reference}

The game targets the mechanical vocabulary of \emph{Wizardry~I}
(grid-based dungeon, turn-based combat, four character classes,
dungeon floors with stairs) and the exploratory sensibility of
\emph{Zork~I} (free-text examination, rich environmental prose,
consequence of curiosity).  Where the two conflict, Wizardry's
determinism wins: rules are authoritative, the engine adjudicates.

\subsection{Design Principles}

\begin{enumerate}[label=\textbf{P\arabic*}.]
  \item \textbf{The engine adjudicates; the LLM narrates.}
    No game outcome is left to probabilistic language model output.
    State transitions, combat resolution, and inventory rules are
    computed deterministically in Go.  The LLM enriches the result
    with prose.
  \item \textbf{Interfaces are the seams.}
    \iface{CommandInterpreter}, \iface{Narrator}, and
    \iface{Renderer} are the only coupling points between the engine
    and external systems.  Each can be replaced independently.
  \item \textbf{Play mode is operational context, not architecture.}
    The same binary supports DM, solo, and headless play.  No code
    paths are conditional on mode; mode selects implementations.
  \item \textbf{Daemon scope stays minimal.}
    The shared game engine daemon handles game logic and push
    routing only.  No LLM calls, no orchestration, no business
    logic outside the game model.
  \item \textbf{Save files are cross-mode.}
    A save created in \mode{dm} mode is loadable in \mode{solo} or
    \mode{headless}.  The \field{play\_mode} field in the save is
    advisory; \cmd{--mode} overrides it at resume time.
  \item \textbf{Party is designed in, not bolted on.}
    \texttt{GameState.Party} is always a slice (length~1 for single
    player).  Combat, movement, and save tooling accept an optional
    \field{character\_id}.  Multi-player (Biff Phase~5) requires no
    engine redesign.
\end{enumerate}

% ===================================================================
\section{Play Modes}
% ===================================================================

Three play modes are first-class citizens of the system.  Each is a
named triple of interface implementations.  The engine always runs as
\cmd{cryptd serve}; there is no embedded deployment.

\begin{center}
\small
\begin{tabular}{lp{3.2cm}p{3.2cm}p{3.2cm}}
\toprule
\textbf{Interface} & \textbf{dm} & \textbf{solo} & \textbf{headless} \\
\midrule
\iface{CommandInterpreter}
  & \texttt{LLMInterpreter}
  & \texttt{SLMInterpreter}
  & \texttt{RulesInterpreter} \\
\iface{Narrator}
  & \texttt{LLMNarrator}
  & \texttt{SLMNarrator}
  & \texttt{TemplateNarrator} \\
\iface{Renderer}
  & \texttt{RPCRenderer} or \texttt{LuxRenderer}
  & \texttt{RPCRenderer} or \texttt{CLIRenderer}
  & \texttt{CLIRenderer} \\
\textbf{Engine deployment}
  & daemon
  & daemon
  & daemon \\
\midrule
\textbf{Invoked by}
  & \cmd{crypt} client \textrightarrow{} \cmd{cryptd serve}
  & \cmd{crypt} client \textrightarrow{} \cmd{cryptd serve}
  & \cmd{cryptd serve -t} \\
\bottomrule
\end{tabular}
\end{center}

\subsection{Inference Tier Chain}

The server auto-selects the best available inference tier at startup.
Each tier provides both a \iface{Command\-Interpreter} and a
\iface{Narrator} implementation.  The chain is:

\begin{enumerate}
  \item \textbf{Large (Claude API)}:
    {\sloppy
    \texttt{LLM\-Interpreter} + \texttt{LLM\-Narrator} via
    OpenAI-compatible \texttt{/v1/chat/completions} endpoint at
    \texttt{api.anthropic.com}.  Activated by
    \texttt{CRYPTD\_API\_KEY} env var or \cmd{--api-key} flag.
    A startup probe validates the key before committing.\par}
  \item \textbf{Medium/Small (SLM)}: \texttt{SLMInterpreter} +
    \texttt{SLMNarrator} via ollama or any OpenAI-compatible server.
    Auto-detected by probing default endpoints.  Preferred models:
    \texttt{gemma3:1b}, \texttt{llama3.2:3b}.
  \item \textbf{Tiny (Rules+Templates)}: \texttt{RulesInterpreter} +
    \texttt{TemplateNarrator}.  Deterministic fallback when no
    inference server is available.  Zero external dependencies.
\end{enumerate}

Both LLM and SLM interpreters use a rules-first strategy:
deterministic commands (``go north'', ``inventory'') are handled by
\texttt{RulesInterpreter} without an inference call.  The model is
only invoked when the rules interpreter returns \texttt{unknown} or
when the action references items/enemies by ambiguous name.

\subsection{Mode Responsibilities}

\begin{center}
\begin{tabular}{lp{2.8cm}p{2.8cm}p{2.8cm}}
\toprule
\textbf{Concern} & \textbf{dm} & \textbf{solo} & \textbf{headless} \\
\midrule
Room narration        & LLM             & SLM             & Template \\
Free-text parsing     & LLM             & SLM             & Rules/regex \\
Scenario generation   & LLM interactive & Pre-made only   & Pre-made only \\
Encounter lore        & LLM             & Pre-made only   & Pre-made only \\
State transitions     & Engine          & Engine          & Engine \\
Combat resolution     & Engine          & Engine          & Engine \\
Fog-of-war, save/load & Engine          & Engine          & Engine \\
Rendering             & Lux or RPC      & RPC or CLI      & CLI \\
\bottomrule
\end{tabular}
\end{center}

% ===================================================================
\section{System Architecture}
% ===================================================================

\subsection{Component Map}

\begin{center}
\begin{tabular}{lp{3cm}p{7cm}}
\toprule
\textbf{Component} & \textbf{Language / Layer} & \textbf{Responsibility} \\
\midrule
Game engine
  & Go (L1)
  & All state transitions; character, map, inventory, combat, save/load.
    Zero LLM dependencies. \\
\iface{CommandInterpreter}
  & Go interface
  & Maps player free-text to a structured \texttt{EngineAction}. \\
\iface{Narrator}
  & Go interface
  & Generates human-readable prose from engine event results. \\
\iface{Renderer}
  & Go interface
  & Presents game state to the player (Lux HUD or ANSI CLI). \\
\cmd{cryptd} server
  & Go binary
  & JSON-RPC 2.0 daemon; exposes direct named methods
    (\texttt{session.init}, \texttt{game.new}, \texttt{game.play},
    \texttt{game.move}, \dots) over Unix socket or TCP.
    Game-as-goroutine architecture: each game is a goroutine that
    exclusively owns its engine and state.  Session registry tracks
    connections; game registry tracks game goroutines.  Session mode
    (normal: interpreter + narrator + RPCRenderer; passthrough:
    direct \texttt{game.*} dispatch) is selected per-session during
    \texttt{session.init}. \\
LLM inference
  & Go (\pkg{internal/\allowbreak inference})
  & OpenAI-compatible HTTP client for
    \texttt{/v1/\allowbreak chat/\allowbreak completions}.
    Used by \texttt{LLM\-Interpreter}, \texttt{LLM\-Narrator},
    \texttt{SLM\-Interpreter}, and \texttt{SLM\-Narrator}.  Supports
    Claude API (with API key auth) and ollama (no auth). \\
\cmd{crypt} client
  & Go binary
  & Thin CLI client; connects to \cmd{cryptd serve}, sends text, displays
    game with readline and HP/MP bars. Auto-starts server. \\
\cmd{crypt-admin}
  & Go binary
  & Scenario authoring; graph-first generation, validation, SQLite
    working format, YAML directory export (DES-027). \\
\cmd{eval-balance}
  & Go binary
  & Monkey test harness; parallel game sessions with weighted-random
    actions, JSON aggregate reports for balance tuning. \\
\pkg{internal/monkeytest}
  & Go package
  & \texttt{MonkeyRenderer}, \texttt{SessionMetrics}, parallel runner.
    Plugs into the game loop via the Renderer interface. \\
\pkg{internal/scengen}
  & Go package
  & Graph types, topology sources, visitors, YAML exporter, SQLite store.
    Used by \cmd{crypt-admin}; SQLite not linked into server or client. \\
SKILL.md
  & Prompt (L4)
  & DM role instructions for Claude; narration, scenario generation,
    lore, examination responses. \\
Lux
  & Python / ImGui
  & Native display surface; receives element trees from
    \texttt{LuxRenderer} via MCP tools. \\
ollama
  & External service
  & Serves SLM for \texttt{SLMInterpreter} and \texttt{SLMNarrator}.
    Auto-detected in interactive \cmd{-t} mode; not required for CI. \\
\bottomrule
\end{tabular}
\end{center}

\subsection{Topology}

The engine always runs as \cmd{cryptd serve}.  There is no embedded
topology.  Two modes determine how clients interact with the server:

\begin{lstlisting}[title={Normal mode (crypt client -> cryptd serve)}]
crypt (thin client)
    |
    | NDJSON over Unix socket or TCP
    v
cryptd serve
    |-- Session Registry (map[sessionID]*Session)
    |-- Game Registry (map[gameID]*Game)
    |
    |-- Connection 1 (goroutine)
    |       `-- Session A -> Game X (goroutine)
    |                           |-- engine (exclusive)
    |                           |-- state (exclusive)
    |                           `-- commands channel
    |-- Connection 2 (goroutine)
    |       `-- Session B -> Game Y (goroutine)
    |
    `-- Shared: interpreter, narrator (stateless)
\end{lstlisting}

Normal mode flow: client sends free text via \texttt{play} JSON-RPC
method.  The server interprets it (via the selected inference tier),
executes the resulting action on the engine, narrates the outcome,
and returns display-ready text.  The game loop runs inside the game
goroutine via \texttt{RunLoop}, driven by an \texttt{RPCRenderer}
that bridges the network connection to the game loop's
\iface{Renderer} interface.

\begin{lstlisting}[title={Passthrough mode (crypt mcp -> cryptd serve)}]
Claude Code session
    |
    | stdio MCP <-> crypt mcp bridge
    |
    | NDJSON over Unix socket or TCP
    v
cryptd serve
    |-- Session/Game registries (same as normal)
    |-- session.init { mode: "passthrough" }
    |-- game.* methods dispatched directly to engine
    `-- Structured JSON result returned (no interpretation/narration)
\end{lstlisting}

Passthrough mode: the client selects mode by sending
\texttt{session.init} with \field{mode}=\texttt{"passthrough"}.  No
interpreter or narrator is involved.  Each \texttt{game.*} request
(\texttt{game.move}, \texttt{game.look}, \texttt{game.attack}, \dots)
is routed through the game goroutine's command channel for
serialized state access and returns a structured JSON result.  The
LLM (Claude) acts as its own interpreter and narrator.  The daemon
serves normal and passthrough sessions simultaneously on the same
socket; mode is per-session, not per-server.

\subsection{Data Flow: Normal Mode}

In normal mode, the game loop runs inside the game goroutine.  The
\texttt{RPCRenderer} bridges the network connection to the loop.

\begin{enumerate}
  \item \textbf{Connection.}
    Client connects and sends \texttt{session.init} with an optional
    \field{session\_id} and optional \field{mode} (omitted or
    \texttt{""} selects normal mode).  Server returns
    \field{session\_id} and \field{has\_game} (true if this session
    already has an active game from a previous connection).  The
    response is a direct JSON result, not a \texttt{ToolResult}
    wrapper.
  \item \textbf{Game start or resume.}
    If \field{has\_game} is false, client calls \texttt{game.new}
    with scenario and character params.  Server creates a game
    goroutine, sends the initial room, then enters the game loop via
    \texttt{Game.RunLoop()}.  If \field{has\_game} is true, the
    server enters \texttt{resume\-Game\-Loop} immediately --- no
    \texttt{game.new} needed.  The current room is re-rendered.
  \item \textbf{Play loop.}
    Client sends \texttt{game.play} requests with free text.  Inside
    the game goroutine: \texttt{RPC\-Renderer.\-Events()} reads
    the text, the interpreter converts it to an
    \texttt{Engine\-Action}, the engine executes it, the narrator
    generates prose, and \texttt{RPC\-Renderer.\-Render()} writes
    a \texttt{Play\-Response} back to the client.
  \item \textbf{Quit or death.}
    The game loop exits.  A final \texttt{PlayResponse} with
    \field{quit} or \field{dead} set to true is sent.  The game
    goroutine remains alive (state preserved for session resume).
\end{enumerate}

The inference tier (LLM, SLM, or rules+templates) is transparent to
this flow.  Solo and DM modes differ only in which interpreter and
narrator implementations are injected at startup.

\subsection{Data Flow: Passthrough Mode}

Passthrough mode has no game loop.  Each \texttt{game.*} JSON-RPC
request is dispatched to the game goroutine via \texttt{Game.Send()}.
The game goroutine executes the method, returns structured JSON,
and the handler writes it back as a direct JSON-RPC result.  Claude
Code (running as the \cmd{crypt mcp} stdio bridge) acts as its own
interpreter and narrator.

% ===================================================================
\subsection{Daemon Concurrency Model}
% ===================================================================

\subsubsection{Game-as-Goroutine}

Each \texttt{Game} is a goroutine that exclusively owns its
\texttt{*engine.Engine} and \texttt{*model.GameState}.  No mutex
protects game state --- serialization is by construction via a
buffered command channel (\texttt{chan Command}, capacity~1).
This eliminates an entire class of concurrency bugs: the game
goroutine is the only code that reads or writes engine state.

\subsubsection{Command Types}

Three command types flow through the channel:

\begin{center}
\begin{tabular}{lp{8cm}}
\toprule
\textbf{Type} & \textbf{Purpose} \\
\midrule
\texttt{tool}
  & Method dispatch: execute a named \texttt{game.*} JSON-RPC method
    against the engine and return the result.  Used by passthrough
    mode for every request and by normal mode for \tool{game.new}. \\
\texttt{run\_loop}
  & Normal mode game loop: the game goroutine runs
    \texttt{game.NewLoop().Run()} internally, blocking until the
    player quits, dies, or disconnects.  The \texttt{RPCRenderer}
    reads/writes the network connection from within the goroutine. \\
\texttt{inspect}
  & Safe state access: runs a caller-provided function inside the
    game goroutine with exclusive access to engine and state.
    Used for test assertions and handler-level state queries. \\
\bottomrule
\end{tabular}
\end{center}

\subsubsection{Session/Game Separation}

A \texttt{Session} is connection identity --- it persists across
reconnects.  A \texttt{Game} is an engine + state goroutine.
Sessions and games are stored in separate registries
(\texttt{map[string]*Session} and \texttt{map[string]*Game}),
protected by a shared \texttt{sync.RWMutex}.

Today the mapping is 1:1 (one session owns one game).  The
architecture supports N:1 for future multiplayer: multiple sessions
can reference the same game ID, with the game goroutine serializing
all their commands.

\subsubsection{Session Resume}

Client sends \texttt{initialize} with its previous
\field{session\_id}.  The server looks up the session in the
registry.  If the session exists and has a \field{gameID}, the
\texttt{InitializeResult} sets \field{has\_game} to true.  The
handler then enters \texttt{resumeGameLoop}, which sends a
\texttt{run\_loop} command with \texttt{SkipInitialRender=false}
--- the current room is re-rendered so the player sees their
position after reconnecting.

\subsubsection{Graceful Shutdown}

\texttt{s.cancel()} cancels the server context.  All active
connections are closed (unblocking any blocked scanners).
\texttt{wg.Wait()} blocks until all connection handler goroutines
exit.  Game goroutines exit via context cancellation (the
\texttt{select} on \texttt{ctx.Done()} in \texttt{Game.run()}).

\subsubsection{Panic Recovery}

Both game goroutines and connection handler goroutines recover from
panics via \texttt{defer recover()}.  A panic in one game or
connection does not crash the daemon or affect other sessions.

% ===================================================================
\section{The Three Engine Interfaces}
\label{sec:interfaces}
% ===================================================================

\subsection{CommandInterpreter}

Maps a player's free-text string to a structured
\texttt{EngineAction} that the engine can execute deterministically.

\begin{lstlisting}[language=Go]
type CommandInterpreter interface {
    Interpret(ctx context.Context, input string,
              state GameState) (EngineAction, error)
}
\end{lstlisting}

\begin{center}
\small
\begin{tabular}{lp{2.5cm}p{7cm}}
\toprule
\textbf{Implementation} & \textbf{Used in} & \textbf{Mechanism} \\
\midrule
\texttt{LLMInterpreter}
  & \mode{dm} (Large tier)
  & Claude API via OpenAI-compatible \texttt{/v1/\allowbreak chat/\allowbreak completions};
    rules-first strategy (exact verbs handled without LLM call;
    LLM invoked only for ambiguous input needing ID resolution). \\
\texttt{SLMInterpreter}
  & \mode{solo} (Medium tier)
  & ollama or any OpenAI-compatible server; uses
    \pkg{internal/\allowbreak inference} client; same rules-first fallback. \\
\texttt{Rules\-Interpreter}
  & \mode{headless} (Tiny tier)
  & Regex/keyword matching over known verb--noun pairs; no model
    required.  Also serves as fallback for both LLM and SLM
    interpreters on inference failure. \\
\bottomrule
\end{tabular}
\end{center}

\subsection{Narrator}

Generates human-readable prose from an \texttt{EngineEvent}: the
result of an action after the engine has resolved it.

\begin{lstlisting}[language=Go]
type Narrator interface {
    Narrate(ctx context.Context, event EngineEvent,
            state GameState) (string, error)
}
\end{lstlisting}

\begin{center}
\small
\begin{tabular}{lp{2.5cm}p{7cm}}
\toprule
\textbf{Implementation} & \textbf{Used in} & \textbf{Mechanism} \\
\midrule
\texttt{LLMNarrator}
  & \mode{dm} (Large tier)
  & Claude generates rich contextual prose for atmospheric events
    (room descriptions, combat starts, deaths, level-ups, item
    examinations).  Tactical events (attack damage, inventory
    changes) use the template fallback.  Selective enrichment:
    LLM for flavour, template for precision. \\
\texttt{SLMNarrator}
  & \mode{solo} (Medium tier)
  & ollama generates 1--3 sentences of atmospheric prose from room
    seed + event type.  Same selective enrichment strategy as LLM. \\
\texttt{Template\-Narrator}
  & \mode{headless} (Tiny tier)
  & Go text/template: ``You move north. You are in the
    \{room.name\}.''  Also serves as fallback for both LLM and SLM
    narrators on inference failure. \\
\bottomrule
\end{tabular}
\end{center}

\subsection{Renderer}

Presents the current \texttt{GameState} to the player after each
turn.  Also receives button interaction events in DM/solo modes.

\begin{lstlisting}[language=Go]
type Renderer interface {
    Render(ctx context.Context, state GameState,
           narration string) error
    Events() <-chan InputEvent  // button presses, free-text
}
\end{lstlisting}

\begin{center}
\small
\begin{tabular}{lp{2.8cm}p{7cm}}
\toprule
\textbf{Implementation} & \textbf{Used in} & \textbf{Mechanism} \\
\midrule
\texttt{RPCRenderer}
  & Normal mode daemon
  & Bridges the game loop to the network connection in normal mode.
    \texttt{Render()} serializes a \texttt{Play\-Response} as NDJSON.
    \texttt{Events()} reads \texttt{play} requests from the client's
    scanner.  Runs inside the game goroutine. \\
\texttt{LuxRenderer}
  & \cmd{-t --lux} mode
  & Calls Lux MCP tools (\tool{show}, \tool{update}, \tool{recv})
    to drive the full HUD.  Currently used in \cmd{-t} testing mode. \\
\texttt{CLIRenderer}
  & \cmd{-t} mode, \mode{headless}
  & ANSI terminal; ASCII map, text HP bar, stdin input loop. \\
\bottomrule
\end{tabular}
\end{center}

% ===================================================================
\section{Game Engine Design}
\label{sec:engine}
% ===================================================================

\subsection{Core Data Model}

\begin{lstlisting}[language=Go]
type GameState struct {
    ScenarioID  string
    PlayMode    string   // advisory; --mode overrides on load
    Timestamp   string
    Character   Character
    Party       []Character  // len 1 for single-player
    Dungeon     DungeonState
    Log         []LogEntry
}

type Character struct {
    Name       string
    Class      CharClass   // Fighter | Mage | Thief | Priest
    Level      int
    HP, MaxHP  int
    MP, MaxMP  int
    Stats      Stats       // STR INT DEX CON WIS CHA
    XP, Gold   int
    Inventory  []Item
    Equipped   Equipment
    Conditions []Condition
}

type DungeonState struct {
    CurrentRoom  string
    Visited      []string
    RoomState    map[string]RoomState
}
\end{lstlisting}

\subsection{Character}

\begin{itemize}
  \item \textbf{Classes}: Fighter, Mage, Thief, Priest.
    Class determines HP dice, MP pool, available spells, and XP
    thresholds.
  \item \textbf{Stats}: STR, INT, DEX, CON, WIS, CHA.
    Bonuses computed by standard table (8--18 range, D\&D-adjacent).
  \item \textbf{Equipment slots}: weapon, armour, ring, amulet.
  \item \textbf{Leveling}: XP thresholds defined per class in
    scenario YAML.  Stat gains on level-up are deterministic.
\end{itemize}

\subsection{Map}

\begin{itemize}
  \item Grid-based dungeon in the Wizardry~I style (20$\times$20
    floor).  Rooms are nodes; corridors are edges.
  \item \textbf{Pre-generated}: the DM creates the room graph before
    play begins, either by writing scenario YAML directly or via the
    \cmd{/crypt:create} skill command.  Procedural generation is
    explicitly not used --- it produces maps that are inconsistent on
    backtracking.
  \item \textbf{Fog of war}: unvisited rooms are opaque on the Lux
    map canvas.  Room content is not revealed until the player enters.
  \item \textbf{Room types}: corridor, chamber, stairway (up/down),
    treasure vault, boss chamber.
  \item \textbf{Door kinds}: open, locked (requires key item),
    hidden (requires search action), one-way, trapped.
\end{itemize}

\subsection{Combat}

Turn-based, deterministic.  No probabilities are left to the LLM.

\begin{itemize}
  \item Initiative roll (DEX modifier + 1d10) determines turn order.
  \item \textbf{Player actions per turn}: Attack, Cast Spell, Use
    Item, Flee (DEX check), Defend (+AC until next turn).
  \item \textbf{Enemy AI}: patterns encoded in scenario YAML.
    Supported patterns: \texttt{aggressive} (always attack),
    \texttt{cautious} (flee below 30\% HP), \texttt{support}
    (buff allies before attacking), \texttt{scripted} (fixed action
    sequence).
  \item Damage formula: weapon dice + STR modifier vs.\ target AC.
  \item Conditions: poisoned, stunned, confused, blessed, hasted.
    Duration in turns; applied/cleared by spells and items.
  \item XP is awarded at \texttt{end\_combat}; loot rolls are
    seeded from scenario YAML drop tables.
\end{itemize}

% ===================================================================
\section{Data Formats}
% ===================================================================

\subsection{Scenario Format (YAML)}

Scenarios are pre-authored YAML files stored in \texttt{scenarios/}.
The engine parses them at startup; the LLM never sees the raw YAML
during play.

\begin{lstlisting}[title={scenarios/classic-fantasy-dungeon.yaml (excerpt)}]
id: classic-fantasy-dungeon
title: "The Grimhold Dungeon"
description: "A classic dungeon crawl..."
version: "1.0"
xp_thresholds:
  fighter: [0, 2000, 4000, 8000, 16000]
  mage:    [0, 2500, 5000, 10000, 20000]

rooms:
  - id: entrance
    name: "Dungeon Entrance"
    description: "Stone stairs descend into darkness."
    exits:
      south: goblin_lair
    items: []
    enemies: []

  - id: goblin_lair
    name: "Goblin Lair"
    description: "The reek of smoke and old blood."
    exits:
      north: entrance
      east:  { room: treasure_vault, door: locked, key: rusty_key }
    enemies:
      - template: goblin
        count: 3
    loot_table:
      - item: rusty_key
        drop_chance: 0.8

enemy_templates:
  goblin:
    hp_dice: "2d6"
    ac: 12
    attack_dice: "1d4"
    xp: 30
    ai: aggressive
\end{lstlisting}

\subsection{Save File Format (JSON)}

Save files live at \texttt{.dungeon/saves/<slot>.json} (gitignored).

\begin{lstlisting}[title={.dungeon/saves/1.json (excerpt)}]
{
  "schema_version": "1.0",
  "play_mode": "dm",
  "scenario": "classic-fantasy-dungeon",
  "timestamp": "2026-03-10T18:00:00Z",
  "party": [
    {
      "id": "hero",
      "name": "Aldric", "class": "fighter", "level": 3,
      "hp": 45, "max_hp": 60, "mp": 0, "max_mp": 0,
      "xp": 1240, "gold": 48,
      "stats": {"str":16,"int":9,"dex":12,"con":14,"wis":10,"cha":11},
      "inventory": [],
      "equipped": {"weapon_id": "rusty_sword"},
      "conditions": []
    }
  ],
  "dungeon": {
    "current_room": "goblin_lair",
    "visited_rooms": ["entrance", "corridor_a", "goblin_lair"],
    "room_state": {
      "goblin_lair":    {"cleared": true, "items": []},
      "treasure_vault": {"cleared": false, "items": ["rusty_key"]}
    },
    "combat": {"active": false}
  },
  "adventure_log": []
}
\end{lstlisting}

\texttt{party} is always present (single-element array for single-player).
\texttt{schema\_version} is a string; \texttt{internal/save} refuses loads
when it does not match \texttt{save.SchemaVersion}.  \texttt{play\_mode}
is advisory metadata captured at save time.

% ===================================================================
\section{JSON-RPC Method Surface (Game Engine)}
% ===================================================================

The engine exposes game actions as direct named methods on JSON-RPC
2.0.  There is no MCP \texttt{tools/call} framing: each method has
its own params shape and returns a direct JSON result object (or a
standard JSON-RPC error).

\subsection{Protocol Methods}

\begin{center}
\begin{tabular}{lp{4cm}p{5cm}}
\toprule
\textbf{Method} & \textbf{Params} & \textbf{Returns} \\
\midrule
\texttt{session.init}
  & \field{session\_id?}, \field{mode?}
    (\texttt{""} or \texttt{"passthrough"})
  & \field{protocolVersion}, \field{serverInfo},
    \field{capabilities}, \field{session\_id},
    \field{has\_game} \\
\texttt{session.quit}
  & ---
  & \field{ok} \\
\texttt{game.new}
  & \field{scenario\_id?}, \field{character\_name?},
    \field{character\_class?}, \field{stats?}
  & Initial room state (normal: \texttt{PlayResponse};
    passthrough: structured room JSON) \\
\texttt{game.list\_scenarios}
  & ---
  & \field{scenarios[]}: \field{id}, \field{title},
    \field{description} \\
\texttt{game.list\_sessions}
  & ---
  & \field{sessions[]}: session + character metadata \\
\texttt{game.play}
  & \field{text}
  & \texttt{PlayResponse}: \field{text}, \field{state?},
    \field{dead?}, \field{quit?}, \field{exits?},
    \field{next\_level\_xp?} \\
\texttt{game.context}
  & ---
  & DM-facing scenario/room context \\
\texttt{game.move}
  & \field{direction}
  & Room transition result \\
\texttt{game.look}
  & ---
  & Current room description \\
\texttt{game.attack}, \texttt{game.defend}, \texttt{game.flee}
  & (action-specific)
  & Combat result \\
\texttt{game.examine}, \texttt{game.inventory},
  \texttt{game.drop}, \texttt{game.equip}, \texttt{game.unequip},
  \texttt{game.use\_item}
  & (action-specific)
  & Item/inventory result \\
\texttt{game.save\_game}, \texttt{game.load\_game}
  & \field{slot?}
  & Slot name + file path (save); room state (load) \\
\bottomrule
\end{tabular}
\end{center}

\subsection{Navigation Methods}

\begin{center}
\begin{tabular}{lp{4cm}p{5cm}}
\toprule
\textbf{Method} & \textbf{Params} & \textbf{Returns} \\
\midrule
\texttt{game.move}
  & \field{direction}
  & Room, exits, items; combat if enemies present \\
\texttt{game.look}
  & ---
  & Current room description, exits, items, combat state \\
\bottomrule
\end{tabular}
\end{center}

\subsection{Inventory Methods}

\begin{center}
\begin{tabular}{lp{4cm}p{5cm}}
\toprule
\textbf{Method} & \textbf{Params} & \textbf{Returns} \\
\midrule
\texttt{game.pick\_up}  & \field{item\_id} & Item details, hero summary \\
\texttt{game.drop}      & \field{item\_id} & Item details, hero summary \\
\texttt{game.equip}     & \field{item\_id} & Item, slot, hero summary \\
\texttt{game.unequip}   & \field{slot}     & Item, slot, hero summary \\
\texttt{game.examine}   & \field{item\_id} & Item details \\
\texttt{game.inventory} & ---              & Items, equipped, weight, capacity \\
\texttt{game.use\_item} & \field{item\_id} & Effect applied, hero summary \\
\bottomrule
\end{tabular}
\end{center}

\subsection{Combat Methods}

\begin{center}
\begin{tabular}{lp{4cm}p{5cm}}
\toprule
\textbf{Method} & \textbf{Params} & \textbf{Returns} \\
\midrule
\texttt{game.attack}
  & \field{target\_id?}
  & Damage, killed, hero summary, enemy turns \\
\texttt{game.defend}
  & ---
  & Defence status, hero summary, enemy turns \\
\texttt{game.flee}
  & ---
  & Success (DEX check), hero summary \\
\texttt{game.cast\_spell}
  & \field{spell\_id}, \field{target\_id?}
  & Spell effect, MP cost, hero summary \\
\bottomrule
\end{tabular}
\end{center}

% ===================================================================
\section{Lux Display Integration}
% ===================================================================

\textbf{Status}: partially implemented.  The \texttt{LuxRenderer}
and Lux element tree builder exist and work in \cmd{cryptd serve -t
--lux} mode.  The full four-panel HUD described below is future work.

\subsection{HUD Layout}

The \texttt{LuxRenderer} drives a four-panel HUD.  Navigation and
combat buttons deliver input events directly to the engine --- no
LLM round-trip is required for button-driven actions.

\begin{lstlisting}[title={Lux HUD layout (schematic)}]
+--------------------------------------------------+
| Menu: [Game] [Character] [Map]                   |
+----------------------------+---------------------+
|                            | HP  [========--] 45 |
|   DUNGEON MAP              | MP  [----------]  0 |
|   (draw canvas,            | Fighter Lv3  XP 1k  |
|    grid rooms,             +---------------------+
|    fog of war,             | EXITS               |
|    visited/unvisited,      |       [N]            |
|    player dot,             |  [W]       [E]      |
|    locked doors)           |       [S]            |
|                            |       [D stairs]     |
+----------------------------+---------------------+
| NARRATION LOG              | COMBAT (when active) |
| (markdown, scrolling)      | Goblin [=====--] 8hp |
|                            | [Attack] [Use Item]  |
|                            | [Defend] [Flee]      |
|                            | Round 2 - Your turn  |
+----------------------------+---------------------+
\end{lstlisting}

\subsection{Update Strategy}

\begin{center}
\begin{tabular}{lp{8cm}}
\toprule
\textbf{Event} & \textbf{Lux operation} \\
\midrule
Room transition (new room entered)
  & \tool{lux.show()} --- replace full narration panel, update map fog \\
HP/MP change during combat
  & \tool{lux.update()} --- patch progress bar elements in-place \\
Narration text appended
  & \tool{lux.update()} --- append to narration log list \\
Combat begins/ends
  & \tool{lux.show()} --- activate or deactivate combat panel \\
Map fog reveal (new room)
  & \tool{lux.update()} --- patch map canvas element \\
Nav button press (player)
  & \tool{lux.recv()} event \textrightarrow{} engine directly \\
\bottomrule
\end{tabular}
\end{center}

\subsection{Lux Tools Used}

The \texttt{LuxRenderer} uses the following Lux MCP tools:

\begin{itemize}
  \item \tool{show(scene\_id, elements, title?)} --- initial HUD render
    and full scene transitions.
  \item \tool{update(scene\_id, patches)} --- incremental HP, MP, fog,
    and log updates within a scene.
  \item \tool{recv(timeout?)} --- poll for button interaction events
    (navigation presses, combat actions).
\end{itemize}

The engine does not use \tool{show\_table}, \tool{show\_dashboard},
or \tool{show\_diagram}; those are data-analysis surfaces.

% ===================================================================
\section{mcp-proxy Integration}
% ===================================================================

\textbf{Status}: future / design-stage.  With M10's concurrent
session support and direct TCP/Unix socket connections, mcp-proxy is
less critical than originally envisioned.  Multiple \cmd{crypt}
clients can connect directly to the daemon.  The proxy pattern is
still planned for Claude Code MCP multiplexing, where stdio
constraints require a shim process.

\subsection{Problem}

Claude Code spawns one MCP server process per session via stdio.
If the DM's Claude session owns the engine process via stdio, a
second player connection is architecturally impossible: there is no
shared process to connect to, and the DM cannot push events to
players.

\subsection{Solution}

mcp-proxy is a near-zero-cost Go shim ($<$10~ms startup,
$<$10~MB RSS) that each Claude session spawns via stdio.  The shim
connects to a pre-running shared daemon over a persistent
bidirectional transport (WebSocket or NDJSON Unix socket) and
relays MCP messages.  At connect time, mcp-proxy resolves the
topmost \cmd{claude} ancestor PID and injects it as a session
identity header.

The daemon uses session identity to:

\begin{itemize}
  \item Route DM-privileged methods only to the DM session.
  \item Push \texttt{tools/list\_changed} notifications to the DM
    when a player takes an action (future; see cryptd-90e.3).
  \item Route DM narration output to all player sessions.
\end{itemize}

\subsection{Daemon Scope Constraint}

The beads project accumulated 70\,000 lines of daemon code before
being deleted.  This daemon must not repeat that mistake.  The
game daemon handles:

\begin{itemize}
  \item Game state (authoritative \texttt{GameState} for the session).
  \item JSON-RPC method dispatch (receive call \textrightarrow{}
    engine method \textrightarrow{} return result).
  \item Push routing (future: \texttt{tools/list\_changed},
    narration broadcast).
\end{itemize}

It does not handle: orchestration logic, plugin lifecycle, or
anything outside the game model.  Note: the daemon does now manage
sessions (session registry, resume) and optionally runs inference
(interpreter + narrator in normal mode).  These are scoped
additions, not the unbounded accretion that killed the beads daemon.

% ===================================================================
\section{Inference Integration}
% ===================================================================

All inference calls use the OpenAI-compatible
\texttt{/v1/\allowbreak chat/\allowbreak completions} wire format
via the \texttt{inference.\allowbreak Client} type in
\pkg{internal/\allowbreak inference}.  The same client code works
for Claude API, ollama, llama.cpp, and any other compatible server.

\subsection{Three Tiers}

\begin{center}
\small
\begin{tabular}{lp{3.2cm}p{6.5cm}}
\toprule
\textbf{Tier} & \textbf{Backend} & \textbf{Configuration} \\
\midrule
\textbf{Large}
  & Claude API (\texttt{api.\allowbreak anthropic.\allowbreak com})
  & Activated by \texttt{CRYPTD\_API\_KEY} env var or
    \cmd{--api-key} flag.  Default model:
    \texttt{claude-sonnet-4-20250514} (configurable via
    \cmd{--model}).  Uses \texttt{Authorization: Bearer <key>}
    header.  A startup probe validates the key with a 5~s timeout
    before committing to this tier.  On probe failure, falls
    through to Medium. \\
\textbf{Medium/\allowbreak Small}
  & ollama or any OpenAI-compatible server
  & Auto-detected by probing default endpoints
    (\texttt{localhost:\allowbreak 11434}).  Preferred models:
    \texttt{gemma3:1b}, \texttt{llama3.2:3b}.  Configurable via
    \cmd{--model}. \\
\textbf{Tiny}
  & Rules + Templates (no inference)
  & Deterministic fallback when no inference server is available.
    \texttt{Rules\-Interpreter} + \texttt{Template\-Narrator}.
    Zero external dependencies. \\
\bottomrule
\end{tabular}
\end{center}

\subsection{API Format}

\begin{lstlisting}[title={Inference request (all tiers)}]
POST /v1/chat/completions
{
  "model": "gemma3:1b",
  "messages": [
    {"role": "system", "content": "..."},
    {"role": "user", "content": "..."}
  ],
  "temperature": 0.0,
  "max_tokens": 150
}
\end{lstlisting}

\begin{center}
\begin{tabular}{ll}
\toprule
\textbf{Property} & \textbf{Value} \\
\midrule
API format         & OpenAI-compatible \texttt{/v1/chat/completions} \\
Default SLM model  & gemma3:1b ($\approx$815~MB) \\
Default LLM model  & claude-sonnet-4-20250514 \\
Fallback           & \texttt{RulesInterpreter} on any inference failure \\
Timeout (LLM)      & 15~s \\
Timeout (SLM)      & 5~s \\
Response format    & JSON for interpreters; plain text for narrators \\
\bottomrule
\end{tabular}
\end{center}

\subsection{Interpreter Strategy}

Both \texttt{LLMInterpreter} and \texttt{SLMInterpreter} use a
rules-first strategy: the \texttt{RulesInterpreter} handles the
input first.  If it produces a definitive action without an ambiguous
ID reference, the result is used directly (no inference call).  The
model is invoked only when rules return \texttt{unknown} or when the
action needs ID resolution against game state context.

\subsection{Narrator Strategy}

Both \texttt{LLMNarrator} and \texttt{SLMNarrator} use selective
enrichment: atmospheric events (room descriptions, combat starts,
deaths, level-ups, item examinations) are sent to the model.
Tactical events (attack results, inventory changes, movement
confirmations) use the \texttt{TemplateNarrator} fallback for
precision and speed.

% ===================================================================
\section{Biff Extension Points}
% ===================================================================

The engine is party-ready from day one.  No redesign is required
when Biff multi-player (Phase~5) is added.

\begin{itemize}
  \item \texttt{GameState.Party} is always a \texttt{[]Character}
    slice; single-player sets \texttt{len(Party) == 1}.
  \item \texttt{game.move}, \texttt{game.attack}, \texttt{game.flee},
    \texttt{game.defend} all accept an optional \field{character\_id};
    omitting it targets the primary character.
  \item Combat turn order already handles multiple actors.
  \item Save files carry a \field{party} array (empty until Biff).
\end{itemize}

When Biff is added:

\begin{itemize}
  \item Each Biff participant controls one character in the party.
  \item \texttt{PostToolUse} hooks notify party members of
    state-change events.
  \item The DM narrates for the full party.
  \item Turn tokens are passed via Biff \cmd{/write} messages.
\end{itemize}

% ===================================================================
\section{Implementation Phases}
% ===================================================================

\begin{center}
\small
\begin{tabular}{lp{7.5cm}l}
\toprule
\textbf{Phase} & \textbf{Deliverable} & \textbf{Status} \\
\midrule
\textbf{M0 — Foundation}
  & Project scaffold, Go module, CI, scenario parser, dice parser,
    save format, data contracts (\texttt{GameState}, \texttt{Character},
    \texttt{EngineAction}, \texttt{EngineEvent}).
  & Complete \\
\midrule
\textbf{M1 — Data Contracts}
  & All model types, interface definitions, scenario YAML schema,
    save/load with forward compatibility.
  & Complete \\
\midrule
\textbf{M2 — Thin E2E Slice}
  & End-to-end architecture validation: engine core (move, look,
    pick\_up), \texttt{RulesInterpreter}, \texttt{TemplateNarrator},
    \texttt{CLIRenderer}, \cmd{cryptd serve -t} mode, e2e test
    framework.
  & Complete \\
\midrule
\textbf{M8 — Server Thin Slice}
  & \cmd{cryptd serve} daemon with JSON-RPC 2.0 over Unix socket
    and TCP; \cmd{crypt} thin client; passthrough and normal modes;
    direct named \texttt{game.*} methods; daemonization; PID file
    management.
  & Complete \\
\midrule
\textbf{M9 — LLM Integration}
  & Three-tier inference chain (Claude API \textrightarrow{} ollama SLM \textrightarrow{}
    rules+templates); \texttt{LLMInterpreter}, \texttt{LLMNarrator},
    \texttt{SLMInterpreter}, \texttt{SLMNarrator}; auto-detection;
    \pkg{internal/inference} OpenAI-compatible client.
  & Complete \\
\midrule
\textbf{M10 — Concurrent Sessions}
  & Game-as-goroutine; session registry; game registry; session
    resume via \field{session\_id} in \texttt{session.init}; concurrent
    connection handling; \texttt{RPCRenderer}; panic recovery;
    graceful shutdown.
  & Complete \\
\midrule
\textbf{Future — DM Mode + Lux HUD}
  & Full Lux HUD (map canvas, HP/MP bars, nav buttons);
    SKILL.md rewritten for DM role;
    \cmd{/crypt:create} --- interactive scenario generation.
  & Planned \\
\midrule
\textbf{Future — Rich World}
  & Shops, traps, cursed items; multiple bundled scenarios;
    DM-generated encounters with backstory and puzzles.
  & Planned \\
\midrule
\textbf{Future — Biff Multi-player}
  & Party management (one character per Biff participant);
    N:1 session-to-game mapping;
    turn tokens via Biff \cmd{/write}.
  & Planned \\
\bottomrule
\end{tabular}
\end{center}

% ===================================================================
\section{Open Questions}
% ===================================================================

\begin{enumerate}[label=\textbf{Q\arabic*}.]
  \item \textbf{mcp-proxy dependency.}
    Still open but less urgent.  With M10's direct TCP/Unix socket
    connections and concurrent session support, multiple \cmd{crypt}
    clients connect directly to the daemon.  The proxy pattern is
    still needed for Claude Code's stdio MCP multiplexing but is no
    longer a blocker for multi-client play.
  \item \textbf{SLM model selection.}
    \emph{Resolved.}  Default: \texttt{gemma3:1b}.  Configurable via
    \cmd{--model} flag.  Auto-detected by probing ollama endpoints.
  \item \textbf{Go MCP SDK.}
    \emph{Resolved.}  Hand-rolled JSON-RPC 2.0 in
    \pkg{internal/\allowbreak daemon} and
    \pkg{internal/\allowbreak protocol}.  Approximately 500~lines;
    no external dependency.
  \item \textbf{Dice notation parser.}
    \emph{Resolved.}  Hand-rolled in \pkg{internal/dice}.  Supports
    \texttt{NdS}, \texttt{NdS+M}, \texttt{NdS-M}.
  \item \textbf{Permadeath default.}
    \emph{Resolved.}  Configurable per scenario in YAML.  Hero death
    ends the game session (\field{dead} flag in \texttt{PlayResponse}).
\end{enumerate}

% ===================================================================
\section{Test and Verification Architecture}
\label{sec:testing}
% ===================================================================

The authoritative testing reference is \texttt{docs/testing.md}.
This section provides a brief summary.

\subsection{Philosophy}

The single most important property of the test suite is that it tells
you immediately and unambiguously whether a change broke something.
Three principles follow from this:

\begin{enumerate}[label=\textbf{P\arabic*}.]
  \item \textbf{Determinism is the invariant.}
    The engine (Section~\ref{sec:engine}) must always produce the
    same output for the same input.  Every test that touches engine
    logic asserts a deterministic outcome.  No probabilistic
    assertions, no ``roughly correct.''
  \item \textbf{Interface seams are injection points.}
    \iface{CommandInterpreter}, \iface{Narrator}, and
    \iface{Renderer} (Section~\ref{sec:interfaces}) are the only
    coupling points to external systems.  All tests above unit level
    use test doubles at these seams.  No real LLM, SLM, or Lux
    instance is ever required to run CI.
  \item \textbf{Headless mode is the CI workhorse.}
    \mode{headless} uses \texttt{RulesInterpreter} +
    \texttt{TemplateNarrator} + \texttt{CLIRenderer} --- all
    zero-dependency implementations.  It runs complete game sessions
    programmatically and is the primary vehicle for integration and
    end-to-end tests.
\end{enumerate}

\subsection{Testing Pyramid}

\begin{center}
\begin{tabular}{llrp{5.5cm}}
\toprule
\textbf{Layer} & \textbf{Trigger} & \textbf{Target} & \textbf{Coverage} \\
\midrule
Unit           & Every commit & $<$5~s   & Pure engine functions, parsers, \texttt{RulesInterpreter}, \texttt{TemplateNarrator} \\
Integration    & Every commit & $<$30~s  & Interface compositions, fake SLM, fake Lux, daemon routing, session resume \\
End-to-End     & Every PR     & $<$2~min & Real \cmd{cryptd} subprocess; MCP wire smoke; save round-trips \\
Acceptance     & Release      & $<$10~min & YAML game scripts through \cmd{cryptd serve -t}; full adventures \\
\bottomrule
\end{tabular}
\end{center}

All four layers require zero external services.  The full unit and
integration pass runs in under 30~seconds.  See
\texttt{docs/testing.md} for the full test architecture: layer
details, test double reference, fixture layout, CI configuration,
and verification gates beyond tests (MCP schema contract, race
detection, save format forward compatibility).

\end{document}
